% eksperymentalne tłumaczenie na włoski

%

\subsection{Nierówność Leibniza} % lang-pl
\subsection{Disuguaglianza di Leibniz} % lang-it

U Zetela \cite[s. 72]{zetel_2020} znajdziemy wniosek z twierdzenia Stewarta (\ref{twierdzenie_stewarta}): % lang-pl
In Zetel \cite[p. 72]{zetel_2020} troviamo una conseguenza del teorema di Stewart (\ref{twierdzenie_stewarta}): % lang-it

\begin{proposition}
    Odległość między środkiem ciężkości i środkiem okręgu opisanego na trójkącie wynosi % lang-pl
    La distanza tra il baricentro e il circocentro di un triangolo è % lang-it
    \begin{equation}
        \sqrt{R^2 - \frac 19 \left(a^2 + b^2 + c^2\right)},
    \end{equation}
\end{proposition}

Wynika z niego: % lang-pl
Ne consegue: % lang-it

\begin{proposition}
[nierówność Leibniza] % lang-pl
[disuguaglianza di Leibniz] % lang-it
    Niech... % lang-pl
    Sia... % lang-it
    \begin{equation}
        a^2 + b^2 + c^2 \le 9R^2,
    \end{equation}
    z równością dla trójkąta równobocznego. % lang-pl
    \index{nierówność!Leibniza} % lang-pl
    con uguaglianza nel caso del triangolo equilatero. % lang-it
    \index{disuguaglianza!di Leibniz} % lang-it
\end{proposition}

Guzicki \cite[s. 362, 363, 371-375]{guzicki_2021} dowodzi tego trygonometrycznie: % lang-pl
Guzicki \cite[p. 362, 363, 371-375]{guzicki_2021} dimostra questo risultato mediante la trigonometria: % lang-it
\begin{equation}
    \sin^2 \alpha + \sin^2 \beta + \sin^2 \gamma = 2 (1 + \cos \alpha \cos \beta \cos \gamma) \le \frac 94.
\end{equation}

%